% BHU PYQ -- Transcribed & Corrected
% Paper: PHIMJ-43 / PHIMJ43: An Introduction to Philosophy of Religion
% Exam: B.A. (Hons.) Semester IV Examination, 2025-26 (NEP)
% Category: Major Core
% Department: Department of Philosophy, Faculty of Arts, Banaras Hindu University

\documentclass[11pt,a4paper]{article}
\usepackage[a4paper,top=2.5cm,bottom=2.5cm,left=2.8cm,right=2.8cm,headheight=14pt]{geometry}
\usepackage{amsmath,amssymb}
\usepackage[T1]{fontenc}
\usepackage[utf8]{inputenc}
\usepackage{lmodern,microtype,fancyhdr,enumitem,hyperref,lastpage,parskip}

\hypersetup{
    colorlinks=true,
    urlcolor=black,
    linkcolor=black,
    pdftitle={PHIMJ43: An Introduction to Philosophy of Religion -- BHU Sem IV 2025-26 (NEP)}
}

\pagestyle{fancy}
\fancyhf{}
\renewcommand{\headrulewidth}{0.4pt}
\renewcommand{\footrulewidth}{0.4pt}
\fancyhead[L]{\small \textit{Banaras Hindu University}}
\fancyhead[R]{\small \textit{PHIMJ-43: Philosophy of Religion (Sem IV 2025-26)}}
\fancyfoot[C]{\small Page \thepage\ of \pageref{LastPage}}

\newlist{parts}{enumerate}{2}
\setlist[parts,1]{label=(\alph*),leftmargin=*,itemsep=4pt,topsep=2pt}

\newcommand{\pts}[1]{\hfill \textit{[#1]}}

\begin{document}

\begin{center}
  {\large\bfseries BANARAS HINDU UNIVERSITY}\\[2pt]
  {\normalsize B.A. (Hons.) Semester IV Examination, 2025-26 (NEP)}\\[6pt]
  \rule{0.8\linewidth}{0.5pt}\\[6pt]
  {\large\bfseries DEPARTMENT OF PHILOSOPHY}\\[4pt]
  {\large\bfseries Paper Code: PHIMJ-43 : An Introduction to Philosophy of Religion}\\[4pt]
  {\normalsize Time: 2:30 Hours \hfill Max. Marks: 70}
\end{center}

\rule{\linewidth}{0.4pt}

\smallskip

\noindent \textit{(Write your Roll No. at the top immediately on the receipt of this question paper)}\\[2pt]
\noindent \textit{Note: Attempt all the three Sections - A, B and C according to the instructions given.}\\[1pt]
\noindent \textit{निर्देशानुसार तीनों खण्डों अ, ब तथा स के उत्तर दें।}

\bigskip

\begin{center}
  {\large\bfseries SECTION A / खण्ड-अ}\\[2pt]
  {\normalsize (Long Answer Type Questions / दीर्घ उत्तरीय प्रश्न)}
\end{center}

\noindent \textit{Note: Attempt any three questions out of the following five questions, each in 300 words. Each question carries 10 marks.}\\[2pt]
\noindent \textit{नोट: निम्नलिखित पाँच प्रश्नों में से किसी भी तीन प्रश्नों के उत्तर दीजिए। प्रत्येक उत्तर 300 शब्दों में लिखें। प्रत्येक प्रश्न 10 अंक का है।}

\bigskip

\noindent \textbf{Question 1.} What is the nature of Philosophy of Religion? How is it different from religion? Discuss.\\[1pt]
धर्म-दर्शन की प्रकृति क्या है? यह धर्म से किस प्रकार भिन्न है? विवेचन कीजिए। \pts{10}

\medskip\rule{\linewidth}{0.2pt}\medskip

\noindent \textbf{Question 2.} Explain the different theories regarding the origin of religion.\\[1pt]
धर्म की उत्पत्ति के विभिन्न सिद्धांतों की व्याख्या कीजिए। \pts{10}

\medskip\rule{\linewidth}{0.2pt}\medskip

\noindent \textbf{Question 3.} Explain the Ontological Argument for the existence of God.\\[1pt]
ईश्वर के अस्तित्व के लिए सत्तात्मक तर्क की व्याख्या कीजिए। \pts{10}

\medskip\rule{\linewidth}{0.2pt}\medskip

\noindent \textbf{Question 4.} What is the problem of religious language? Explain its philosophical significance.\\[1pt]
धार्मिक भाषा की समस्या क्या है? इसके दार्शनिक महत्व की व्याख्या कीजिए। \pts{10}

\medskip\rule{\linewidth}{0.2pt}\medskip

\noindent \textbf{Question 5.} What is the problem of evil? Why is it a challenge for theism? Discuss.\\[1pt]
अशुभ की समस्या क्या है? यह ईश्वरवाद के लिए चुनौती क्यों है? विवेचन कीजिए। \pts{10}

\bigskip
\rule{\linewidth}{0.2pt}
\bigskip

\begin{center}
  {\large\bfseries SECTION B / खण्ड-ब}\\[2pt]
  {\normalsize (Short Answer Type Questions / लघु उत्तरीय प्रश्न)}
\end{center}

\noindent \textit{Note: Attempt any four questions out of the following six questions, each in 150 words. Each question carries 5 marks.}\\[2pt]
\noindent \textit{नोट: निम्नलिखित छह प्रश्नों में से किसी भी चार प्रश्नों के उत्तर दीजिए। प्रत्येक उत्तर 150 शब्दों में लिखें। प्रत्येक प्रश्न 5 अंक का है।}

\bigskip

\noindent \textbf{Question 6.} What is the relationship between religion and philosophy?\\[1pt]
धर्म और दर्शन का क्या संबंध है? \pts{5}

\medskip\rule{\linewidth}{0.2pt}\medskip

\noindent \textbf{Question 7.} Is religion based on fear or faith? Discuss.\\[1pt]
क्या धर्म भय पर आधारित है या आस्था पर? विवेचन कीजिए। \pts{5}

\medskip\rule{\linewidth}{0.2pt}\medskip

\noindent \textbf{Question 8.} What is the relation between religion and morality?\\[1pt]
धर्म और नैतिकता का क्या संबंध है? \pts{5}

\medskip\rule{\linewidth}{0.2pt}\medskip

\noindent \textbf{Question 9.} What is the Moral Argument?\\[1pt]
नैतिक तर्क क्या है? \pts{5}

\medskip\rule{\linewidth}{0.2pt}\medskip

\noindent \textbf{Question 10.} What is the analogical theory?\\[1pt]
सादृश्यात्मक सिद्धांत क्या है? \pts{5}

\medskip\rule{\linewidth}{0.2pt}\medskip

\noindent \textbf{Question 11.} Can religion solve the problem of evil?\\[1pt]
क्या धर्म अशुभ की समस्या का समाधान कर सकता है? \pts{5}

\bigskip
\rule{\linewidth}{0.2pt}
\bigskip

\begin{center}
  {\large\bfseries SECTION C / खण्ड-स}\\[2pt]
  {\normalsize (Very Short Answer Type Questions / अति लघु उत्तरीय प्रश्न)}
\end{center}

\noindent \textit{Note: Answer all ten parts of the following question in not more than two sentences. Each question carries 2 marks.}\\[2pt]
\noindent \textit{नोट: अधोलिखित प्रश्नों के सभी दस भागों के उत्तर अधिकतम दो वाक्यों में दीजिए। प्रत्येक प्रश्न 2 अंक का है।}

\bigskip

\noindent \textbf{Question 12.}
\begin{parts}
  \item Who gave the Ontological Argument?\\[1pt]
        सत्तात्मक तर्क किसने दिया? \pts{2}

  \item Who gave the symbolic theory of religious language?\\[1pt]
        धार्मिक भाषा का प्रतीकात्मक सिद्धांत किसने दिया? \pts{2}

  \item Who explained religion through psychology?\\[1pt]
        किसने मनोविज्ञान के माध्यम से धर्म की व्याख्या की? \pts{2}

  \item What is natural evil?\\[1pt]
        प्राकृतिक अशुभ क्या है? \pts{2}

  \item What is Pantheism?\\[1pt]
        सर्वेश्वरवाद क्या है? \pts{2}

  \item Who introduced the term ``Henotheism''?\\[1pt]
        ``हिनोथिस्म'' शब्द को सर्वप्रथम किसने प्रतिपादित किया? \pts{2}

  \item What is theology?\\[1pt]
        धर्मशास्त्र क्या है? \pts{2}

  \item What are the practical aspects of religion?\\[1pt]
        धर्म के व्यावहारिक पक्ष क्या हैं? \pts{2}

  \item According to whom, ``God is an article of faith''?\\[1pt]
        ``ईश्वर आस्था का विषय है,'' यह किसके अनुसार है? \pts{2}

  \item Which religion follows the Holy Bible?\\[1pt]
        कौन सा धर्म पवित्र बाइबिल का अनुसरण करता है? \pts{2}
\end{parts}

\vfill
\rule{\linewidth}{0.4pt}

\begin{center}\small
\href{https://bhu-exam.vercel.app/}{Click here to visit our website}\\[4pt]
\href{https://me-aryan.vercel.app/}{Click here to visit our contributor}
\end{center}

\end{document}
